\section{Conclusions}
Stationary voltage inversion reconstructs flux but does not enforce magnetic
symmetry. Under conservative, quasi-static, reflection-symmetric magnetics,
an even q-current co-energy implies d-even and q-odd flux and reciprocal
differential cross-inductances. Saturation and cross-saturation change
magnitudes and local gains while preserving that ideal parity.

A resistance error adds $\resistanceerror\rotationgenerator\current/\omegae$. A rotor-angle error
rotates both the current argument and the resulting flux. Their combined
local fingerprint occupies the d-odd and q-even channels, with angle sensitivity
$\Ldiff\rotationgenerator\current-\rotationgenerator\flux$. This expression requires only first
flux derivatives and explains the PM projection in q flux and the isotropic
cancellation in d flux.

Both flux components, diverse current points, and controlled speed variation
can separate the two sensitivities, provided the stacked model has rank two
and adequate scaled conditioning. Neither rank in a restricted model nor a
good residual fit removes exact voltage-error confounding or real physical
symmetry breaking. Synthetic voltage reconstruction verifies signs, limiting
cases, parameter recovery, and finite-angle discrepancy without substituting
for measurement validation.

The companion's symmetric co-energy framework retains nonlinear magnetics
while preventing forbidden parity components from becoming legitimate fitted
characteristics. Its diagnostic residuals should support proposed corrections
with preserved originals and explicit application. The next practical step is
controlled pairing of real MeasEval data with measured thermal conditions,
terminal-voltage uncertainty, and an independently checked magnetic reference.
